\documentclass[a4paper]{report}

%====================== PACKAGES ======================

\usepackage[french]{babel}
\usepackage[utf8x]{inputenc}
\usepackage{dirtree}
%pour gérer les positionnement d'images
\usepackage{float}
\usepackage{amsmath}
\usepackage{graphicx}
\usepackage[colorinlistoftodos]{todonotes}
\usepackage{url}
%pour les informations sur un document compilé en PDF et les liens externes / internes
\usepackage{hyperref}
%pour la mise en page des tableaux
\usepackage{array}
\usepackage{tabularx}
\usepackage{graphicx}
%pour utiliser \floatbarrier
%\usepackage{placeins}
%\usepackage{floatrow}
\usepackage{xr}
%espacement entre les lignes
\usepackage{setspace}
%modifier la mise en page de l'abstract
\usepackage{abstract}
%police et mise en page (marges) du document
\usepackage[T1]{fontenc}
\usepackage[top=2cm, bottom=2cm, left=2cm, right=2cm]{geometry}
%Pour les galerie d'images
\usepackage{subfig}

%====================== INFORMATION ET REGLES ======================

%rajouter les numérotation pour les \paragraphe et \subparagraphe
\setcounter{secnumdepth}{4}
\setcounter{tocdepth}{4}

\hypersetup{							% Information sur le document
pdfauthor = {Premier Auteur,
			Deuxième Auteur,
			Troisième Auteur,
    		Quatrième Auteur},			% Auteurs
pdftitle = {Nom du Projet -
			Sujet du Projet},			% Titre du document
pdfsubject = {Mémoire de Projet},		% Sujet
pdfkeywords = {Tag1, Tag2, Tag3, ...},	% Mots-clefs
pdfstartview={FitH}}					% ajuste la page à la largueur de l'écran
%pdfcreator = {MikTeX},% Logiciel qui a crée le document
%pdfproducer = {}} % Société avec produit le logiciel

%======================== DEBUT DU DOCUMENT ========================

\begin{document}

%régler l'espacement entre les lignes
\newcommand{\HRule}{\rule{\linewidth}{0.5mm}}

%page de garde
\begin{titlepage}
\begin{center}

% Upper part of the page. The '~' is needed because only works if a paragraph has started.
\includegraphics[width=0.35\textwidth]{./assets/logo}~\\[1cm]

\textsc{\LARGE Université Côte D'Azur}\\[1.5cm]

\textsc{\Large }\\[0.5cm]

% Title
\HRule \\[0.4cm]

{\huge \bfseries Takenoko\\
https://github.com/Mahe-Thomas/takenoko \\[0.4cm] }

\HRule \\[1.5cm]

% Author and supervisor
\begin{minipage}{0.4\textwidth}
\begin{flushleft} \large
\emph{Auteur:}\\
Jeremy \textsc{Bonsaudo}\\
Matthias \textsc{Percelay}\\
Brandon \textsc{Fontany--Legall}\\
Thomas \textsc{Mahe}
\end{flushleft}
\end{minipage}
\begin{minipage}{0.4\textwidth}
\begin{flushright} \large
\emph{} \\
 \textsc{}\\
\emph{} \\
 \textsc{}
\end{flushright}
\end{minipage}

\vfill

% Bottom of the page
{\large \today}

\end{center}
\end{titlepage}

%page blanche
%\newpage
%~
%ne pas numéroter cette page
%\thispagestyle{empty}
%\newpage

%\input{./abstract.tex}

\tableofcontents
\thispagestyle{empty}
\setcounter{page}{0}
%ne pas numéroter le sommaire

%\newpage

%espacement entre les lignes d'un tableau
\renewcommand{\arraystretch}{1.5}

%====================== INCLUSION DES PARTIES ======================

~
\thispagestyle{empty}
%recommencer la numérotation des pages à "1"
\setcounter{page}{0}
%\newpage

\externaldocument{./string}
\externaldocument{./conception}
\chapter{Présentation du projet}


\section{Sujet}
Le projet consiste à réaliser en Java une version numérique du jeu Takenoko créé par Antoine Bauza. Version textuelle n'étant pas destinée à être jouée par des êtres humains mais par des robots jouant de façon autonome. Dans le Takenoko nous allons endosser le rôle d'un jardinier japonais affairé à répondre au mieux aux requêtes du vénérable empereur. Pour ce faire, nous aurons, durant notre tour de jeu, deux choses à faire : Regarder quel temps il fait et effectuer deux actions parmi les possibilités suivantes : agrandir le jardin, irriguer, bouger le jardinier, bouger le panda. 

\section{Avancement du projet}
L'intégralité des fonctionnalités du jeu est implémentée à l'exception du dé météo. Bien que les joueurs le lancent au début de leur tour, cela n'influe pas sur le dérouler du tour.\\
Cependant une règle n'est pas respectée par les robots, ils peuvent effectuer plus de deux actions par tours.
Nous estimons que deux itérations supplémentaires auraient été nécessaires pour terminer le moteur selon les règles du jeu.
\section{Problématique soulevée}

\subsection{Détecter et réduire la dette technique}
L'organisation du code a été pour nous un assez grand problème notamment à cause de l'apparition de "God class"\footnotemark\\ comme la classe de Jeu qui concentre un trop grand nombre de fonctions. Nous avons donc fait plusieurs réaménagements \textit{voir paragraphe \hyperref[conception]{Conception}} \\ 

À l'inverse des "god classe", la méthode agile nous impose de concevoir le projet comme un ensemble de petites briques. Durant la vie du projet il peut s'avérer que plusieurs briques ont des besoins fonctionnels similaires, l'ajout de code générique devient alors nécessaire pour répondre aux bonnes pratiques de conception orientés objets. Ces phases d'ajout de généricité n'ajoutent aucune plus-value au projet d'un point de vue client mais peux être vues comme un investissement pour le développement futur.

%Contenu de la note précédemment marquée avec \footnotemark
\footnotetext{La "god Class" est une classe comportant un très grand nombre de fonctions qui pourraient être dispatchées dans d'autres classes}

\subsection{L'injection de dépendance}

Du fait de la nouveauté de l'outil et de l'avancée du projet, nous avons rencontré plusieurs problèmes lors de l'injection de dépendance.\\
Bien que la détection des classes nécessitant des injections de dépendance soit "facile", l'ajout d'IoC sur une classe à un impact fort sur les classes l'utilisant. Cela implique de réviser l'ensemble des tests et d'adapter le reste du projet. \textit{Voir paragraphe \hyperref[spring]{Spring}}


\chapter{Organisation du code }

\section{Packages}

\begin{minipage}{0.45\textwidth}
\dirtree{%
.1 main.
.2 resources.
.2 takenoko.
.3 Deque.
.3 Irrigation.
.3 Joueur.
.4 Strategie.
.5 StrategieAction.
.5 StrategieAmenagement.
.5 StrategieCoord.
.5 StrategieIrrig.
.5 StrategieJardinier.
.5 StrategiePanda.
.3 Objectives.
.4 Amenagement.
.4 Patterns.
.3 Plot.
.3 Properties.
.3 Util.
.4 Comparators.
.4 Exceptions.
}
\end{minipage}
\begin{minipage}{0.45\textwidth}
  Après de multiples ré-aménagements, nous sommes arrivé à cette structuration des packages.\\
  
\end{minipage}

\section{Interfaces}
Les principales Interfaces sont celle des Stratégies, chaque stratégie que ce soit StratégieAction, StratégieAmenagement, etc... ont chacune leur Interfaces car plusieurs implémentations avec des idées différentes sont présentes. Toutes ces interfaces permettant l'implémentation d'une stratégie complexe composée d'une stratégie de chaque catégorie. 

\section{Héritage}
Nous utilisons l'héritage pour les éléments ayant comme dit précédement des besoins fonctionnels similaires, nous avons donc du code générique pour les pioches et les cartes objectifs. Les comportements du type pioche/validation d'une carte étant commun à l'essemble des cartes, il est important que les classes spécifique hérite d'une classe contenant les comportements primaires de ces éléments.

\chapter{Conception}
\label{conception}
Intro\footnotemark\\
%note en bas de page

\section{Patron de conception}

%Contenu de la note précédemment marquée avec \footnotemark
\footnotetext{Note bas de page "intro"}

Singleton : Utilisé pour notre Logger, un unique objet est instancié lors du premier appel et ce pendant toute la durée de vie du programme.\\

Prototype : Utilisé pour la quasi totalité des beans Spring. Permet d'instancier plusieurs instances d'objets selon une unique spécification. Très utile pour générer les decks.\\
%retour à la ligne (alinea)

Strategy : Utilisé pour les stratégies de nos robots. Permet d'encapsuler plusieurs stratégies afin d'avoir une stratégie mère se chargeant de produire le résultat espéré selon une action défini.

%\newpage



Bla\footnotemark\\

%Detail :
Bla(cf. ref. \cite{cite6}).
%citation référencé dans le document "bibliographie.bib" inclus à la fin du document

\footnotetext{Note bas de page "bla"}

\chapter{Spring}

Intro\footnotemark\\
%note en bas de page

\section{Notre Spring}

Bla(cf. fig. 1.1)\\



%Contenu de la note précédemment marquée avec \footnotemark
\footnotetext{Note bas de page "intro"}

Bla
%retour à la ligne (alinea)

Bla\\
%saut de paragraphe

Bla

%\newpage



Bla\footnotemark\\

%Detail :
Bla(cf. ref. \cite{cite6}).
%citation référencé dans le document "bibliographie.bib" inclus à la fin du document

\footnotetext{Note bas de page "bla"}

\chapter{Bilan}

Projet dans l'ensemble plutôt about même si il  reste encore pas mal de travail notamment avec l'ajout du dès météo qui est planifié pour une version prochaine. Après pas mal de restructuration du projet les package sont plus ou moins cohérent de notre point de vu mais en moins les stratégies restent difficiles à prendre en main pour un nouveau développeur arrivant sur le projet. De plus, les patrons de conception nous ont été d'une grande utilité notamment pour les stratégies et la conception d'une stratégie complexe dérivant d'une multitude de petites stratégies spécialisées. Pour finir, nous nous sommes heurté à la difficulté de Spring un peux tardivement ce qui nous obligea à modifier énormément de classes dans notre projet.


Pour conclure, le passage de Spring est un passage obliger au vu de la composition des classes notamment la Classe Game qui connect "les fils" entre les différentes fonctionnalités de l'application. Cependant, le passage à spring devrait dans l'idéale se faire dès le début du projet.

%\input{./annexes.tex}

%\newpage

%récupérer les citation avec "/footnotemark"
\nocite{*}

%choix du style de la biblio
%\bibliographystyle{plain}
%inclusion de la biblio
%\bibliography{bibliographie.bib}
%voir wiki pour plus d'information sur la syntaxe des entrées d'une bibliographie

\end{document}